Laying out the PCB components was a difficult problem for v1.0 where I wanted to keep the board as small as possible to emulate how large the flight-model would be. Comparing \autoref{fig:v1-connectors} to early drafts of the PCB like \autoref{fig:v1-init-layout}, we can see that component layout changed quite a bit, especially with the connectors. The main focus of the initial layout was to place components as close to their respective STM32 pins as possible, for example, keeping the trace length between the STM and C2I connector was a must. 
\nl
The components that took priority in the layout were the: STM32, C2I connector, HSE, LSE, and QSPI. I kept the STM32 as close to the middle as possible, angled on a 45° angle to open the pins up to the edges of the board more (see \href{https://www.youtube.com/watch?v=D6dlfS53Tx0&t=473s}{this YouTube} video for more about the 45° angle). The LSE and HSE were placed as close as possible to the STM32. Placing the 32-pin C2I connector was difficult since the DCMI pins on the STM32 are spread across the entire board, but it was placed closest to the PCLK pin since the signal integrity of that trace superceeded all others as mentioned by \href{https://www.reddit.com/r/embedded/comments/1ozbddc/for_dcmi_pcb_routing_using_ov2640_camera_on_stm32/}{this Reddit post}.
\nl
Some components like the USB-C connector, EPS connector, and regulation circuits needed to be grouped together to minimise how much space is used. Moreover, the trigger button needed to be placed on the edge of the PCB in an area that is easily accessible, and have a clearance of more than 20 mm so a finger can actually push the button - if the button was landscape and squished between the microSD and STLINK/V2 headers, I would not be able to fit my finger between the headers to press the button (I would not be able to press it from above when the microSD module was plugged in). Although the USB data traces travel quite far, and placing the voltage regulation circuit on the top-half would reduce the trace length, it is vital to note that the USB COM port functionality is not part of the system requirements, so the signal integrity of all other traces are a priotity.
\nl
With the selection, STLINK/V2, and microSD headers removed, as well as the USB-C connector, future versions of the PCB should be much easier to layout since most of the issues came from making space on the edge of the PCB.

\begin{indentedfigure}
    \centering
    \includegraphics[width=0.8\textwidth]{Images/v1/PCB/Connectors.png}
    \caption{C2Sv1.0 connection interfaces.}
    \label{fig:v1-connectors}
\end{indentedfigure}